\subsection{THE AXIOM OF THE SET OF TWO ELEMENTS}

\[
(\forall x) (\forall y) \operatorname{Coll} _ {z} (z = x \text {   or   } z = y).
\]

This axiom says that if \(x\) and \(y\) are objects, then there is a set whose only elements are \(x\) and \(y\) .

DEFINITION 2. The set \(\mathfrak{E}_z\) ( \(z = x\) or \(z = y\) ), whose only elements are \(x\) and \(y\) , is denoted by \(\{x, y\}\) .

The relation \(z \in \{x, y\}\) is therefore equivalent to `` \(z = x\) or \(z = y\) ''; it follows from C50 that \(\{y, x\} = \{x, y\}\) .

Let \(R\{z\}\) be a relation and let \(x, y\) be letters distinct from \(z\) . From the criteria C32, C33 (Chapter I, §4, no. 3), and C43 (Chapter I, §5, no. 1) it follows easily that the relation \((\exists z)((z \in \{x, y\})\) and \(R\{z\})\) is equivalent to `` \(R\{x\}\) or \(R\{y\}\) ''; consequently the relation \((\forall z)((z \in \{x, y\}) \Longrightarrow R\{z\})\) is equivalent to `` \(R\{x\}\) and \(R\{y\}\) ''.

The set \(\{x, x\}\) , which is denoted simply by \(\{x\}\) , is called the set whose only element is x; the relation \(z \in \{x\}\) is equivalent to z = x, and the relation \(x \in X\) is equivalent to \(\{x\} \subset X\) .

